\documentclass[10pt,a4paper]{amsart}
\usepackage[margin=19mm]{geometry}
\usepackage{amsmath,amssymb,mathtools}
\usepackage[scr=rsfs,cal=euler]{mathalfa}
\usepackage{xfrac}
\usepackage[unicode,hidelinks,bookmarks=false]{hyperref}
\newcommand{\ie}{i.e.\ }
\newcommand{\cf}{cf.\ }
\newcommand{\resp}{respectively\ }
\newcommand{\Subsection}{\subsection}
\providecommand{\nobreakdash}{\nobreak\hbox{-}}
\setlength{\emergencystretch}{2em}
\multlinegap0pt
\usepackage{bm}
\usepackage{bbm}
\usepackage{dsfont}
\usepackage{xspace}
\usepackage{cancel}
\usepackage{mathtools}
\usepackage[shortlabels]{enumitem}
\usepackage{amsthm}
\makeatletter
\@ifundefined{theorem}{\newtheorem{theorem}{Theorem}}{}
\@ifundefined{lemma}{\newtheorem{lemma}[theorem]{Lemma}}{}
\makeatother
\DeclareMathOperator{\conv}{conv}
\DeclareMathOperator{\pos}{pos}
\DeclareMathOperator{\relint}{relint}
\newcommand{\F}{\mathbb F_2}
\title[A Chain-Level Borsuk--Ulam Obstruction Proof of Norine's Ant]{A Chain-Level Borsuk--Ulam Obstruction Proof of Norine's Antipodal-Coloring Conjecture}
\author{Hehui Wu, Ningyuan Yang}
\date{Source: arXiv:2607.19276v1 (CC BY 4.0)}
\begin{document}
\maketitle

\noindent\emph{Source and licence.} A Chain-Level Borsuk--Ulam Obstruction Proof of Norine's Antipodal-Coloring Conjecture. Version arXiv:2607.19276v1 (CC BY 4.0), distributed under the Creative Commons Attribution 4.0 International licence, as recorded in the arXiv licence metadata for this exact version and shown on the abstract page https://arxiv.org/abs/2607.19276v1. Full rights notice: licenses/arxiv-2607.19276-CC-BY-4.0.txt.\par\smallskip

\noindent\textit{Editorial scope.} Counted unit: the Lemma whose source label is \texttt{lem:degenerate-radial-cancellation} in the article (source0.tex lines 1324--1336, source label \texttt{lem:degenerate-radial-cancellation}), delivered together with the article's own complete original proof of it (source0.tex lines 1338--1423, one \texttt{proof} environment of 86 lines). No mathematics is written, repaired or re-derived: statement and proof are the article's own text line for line. Required context retained uncounted: lem:spherical-affine-model (lines 526--545); lem:refinement-membership (lines 686--709); lem:positive-cone-section (lines 1161--1172); eq:low-dimensional-facet-image (lines 1223--1227); lem:generic-degenerate-fiber (lines 1232--1244). Every definition, display and label of those lines is retained unchanged. Omitted from this package and not redistributed: 1--525, 546--685, 710--1160, 1173--1222, 1228--1231, 1245--1323, 1337--1337, 1424--1801. Nothing inside a retained excerpt is omitted or altered: every retained body is delivered exactly as the article's lines are written, so no omission receipt of any kind accompanies this package. Citations: the retained excerpts carry no citation command at all, so no bibliography entry of the article is transcribed and no reference is redistributed. Reference adaptation: none was needed and none was made; every reference inside the delivered unit resolves inside it, so no unresolved reference or citation is produced. Printed slips kept literal: no printed word, symbol, hypothesis or number of the counted statement or of its proof was altered. Layout and class: the article's own class is \texttt{amsart} with packages including geometry, fontenc, lmodern, microtype, amsmath, amssymb, amsthm, mathtools, enumitem, xcolor, hyperref; the wrapper is the lane's pinned \texttt{amsart} editorial wrapper at 10pt on A4, which restates the theorem-like environments the retained text needs and adds the article's own macro definitions that the retained text needs (\texttt{\textbackslash F}, \texttt{\textbackslash conv}, \texttt{\textbackslash pos}, \texttt{\textbackslash relint}), with extra wrapper packages (bm, bbm, dsfont, xspace, cancel, mathtools, enumitem). The delivered numbering is the unit's own. Assets: no retained span contains a figure environment, a graphics-inclusion command, a PDF-page inclusion, a table or a TikZ picture; no image file is copied into this package, the package declares no asset dependency and no image file is redistributed with it. Licence: version arXiv:2607.19276v1 of this article is distributed under the Creative Commons Attribution 4.0 International licence, as recorded in the arXiv licence metadata for this exact version and shown on the abstract page https://arxiv.org/abs/2607.19276v1. A full rights notice is delivered at licenses/arxiv-2607.19276-CC-BY-4.0.txt. Characters: every retained excerpt is pure ASCII, so the delivered engine, XeLaTeX with its default Latin Modern text font, reports no missing character.
\begin{lemma}
\label{lem:spherical-affine-model}
Let \(0\ne\mathcal C\subseteq W_K\) be a pointed polyhedral cone,
and let \(P,\ell,H_\ell,P_\ell\) be the associated objects defined
above.
Then \(P_\ell\) is a compact convex polytope of dimension
\(\dim\mathcal C-1\), and
\[
 \pi_\ell:P\longrightarrow P_\ell,
 \quad x\longmapsto\frac{x}{\ell(x)},
 \qquad
 \kappa_\ell:P_\ell\longrightarrow P,
 \quad y\longmapsto\frac{y}{\|y\|},
\]
are inverse homeomorphisms.  They induce a dimension-preserving
isomorphism of face posets, carry relative interiors to relative
interiors, and identify
face-to-face subdivisions of \(P\) with ordinary face-to-face
polytopal subdivisions of \(P_\ell\).
\end{lemma}

\begin{lemma}
\label{lem:refinement-membership}
Let \(P\) be a spherical \(j\)-polytope subdivided by a spherical
polyhedral complex \(\mathcal K\).

\begin{enumerate}[label=\textup{(\roman*)},leftmargin=*]
\item If \(Q\in\mathcal K\) has dimension \(j\) and
\(\relint Q\cap\relint P\ne\varnothing\), then \(Q\subseteq P\).

\item Suppose that a spherical polyhedral complex \(\mathcal L\)
refines \(\mathcal K\), meaning that it subdivides every member of
\(\mathcal K\).  If \(H\in\mathcal L\) and
\(Q\in\mathcal K\) have dimension \(j\), with
\[
 H\subseteq Q,
 \qquad
 \relint H\subseteq\relint Q,
\]
then
\[
 H\subseteq P\quad\Longleftrightarrow\quad Q\subseteq P.
\]
\end{enumerate}
\end{lemma}

\begin{lemma}
\label{lem:positive-cone-section}
For every nonempty \(I\subseteq\{0,\ldots,j\}\),
\[
 H_\ell\cap\pos\{v_i:i\in I\}=\conv\{w_i:i\in I\},
 \qquad
 \kappa_\ell\bigl(\conv\{w_i:i\in I\}\bigr)
 =S(W_K)\cap\pos\{v_i:i\in I\}.
\]
Moreover, \(\{v_i:i\in I\}\) is linearly independent if and only if
\(\{w_i:i\in I\}\) is affinely independent.
\end{lemma}

\begin{equation}\label{eq:low-dimensional-facet-image}
 \dim F_i\leq j-2
 \quad\Longrightarrow\quad
 F_i\subseteq E.
\end{equation}

\begin{lemma}
\label{lem:generic-degenerate-fiber}
For every
\[
 y\in\operatorname{aff}\{w_0,\ldots,w_j\}\setminus E,
\]
exactly two of the sets \(F_i\) contain \(y\) when
\(y\in f(\Delta^j)\), and none contains \(y\) otherwise.  Consequently,
\[
 \#\{i:y\in F_i,\ \dim F_i=j-1\}=0
 \qquad\text{in }\F.
\]
\end{lemma}

\begin{lemma}
\label{lem:degenerate-radial-cancellation}
Let \(j\geq1\).  Suppose that
\(v_0,\ldots,v_j\in W_K\setminus\{0\}\) are linearly dependent and
\[
 0\notin\conv\{v_0,\ldots,v_j\}.
\]
Then
\[
 \sum_{i=0}^j
 \mathcal R[v_0,\ldots,\widehat{v_i},\ldots,v_j]=0.
\]
\end{lemma}

\begin{proof}
The cone
\[
 \mathcal C=\pos\{v_0,\ldots,v_j\}
\]
is pointed by the convex-hull hypothesis.  Let
\(\ell,\kappa_\ell\) be defined as in
Lemma~\ref{lem:spherical-affine-model}.  Put
\[
 w_i=\frac{v_i}{\ell(v_i)}.
\]
Lemma~\ref{lem:positive-cone-section} shows that
\(\{w_0,\ldots,w_j\}\) is affinely dependent.  If
\[
 \dim\conv\{w_0,\ldots,w_j\}\leq j-2,
\]
then every \(j\)-element subfamily of the \(w_i\) is affinely
dependent.  The same lemma shows that every radial chain in the
asserted sum is zero.

It remains to consider
\(\dim\conv\{w_0,\ldots,w_j\}=j-1\).  Use the notation
\(f,\Delta^j,F_i,E\) from
Lemma~\ref{lem:generic-degenerate-fiber}, and define free chains
\[
 \widehat P_i=
 \begin{cases}
  \langle\kappa_\ell(F_i)\rangle,&\dim F_i=j-1,\\
  0,&\dim F_i\leq j-2,
 \end{cases}
 \qquad
 P_i=\varpi_{j-1}(\widehat P_i).
\]
By Lemma~\ref{lem:positive-cone-section},
\[
 P_i=\mathcal R[v_0,\ldots,\widehat{v_i},\ldots,v_j].
\]

Choose a common refinement \(\mathcal K\) of the spherical polytopes
occurring in the nonzero \(\widehat P_i\).  Let
\(Q\in\mathcal K\) have dimension \(j-1\).  If \(Q\) is contained in
none of these polytopes, then
\(\mu_{\sum_i\widehat P_i}(Q)=0\).  We may therefore assume that
\(Q\) is contained in at least one of them.  Then
\(\kappa_\ell^{-1}(\relint Q)\) lies in
\(\operatorname{aff}\{w_0,\ldots,w_j\}\).  This relative interior has
dimension \(j-1\), whereas \(E\) is a finite union of polytopes of
dimension at most \(j-2\).  Choose
\[
 y\in\kappa_\ell^{-1}(\relint Q)\setminus E,
 \qquad
 z=\kappa_\ell(y)\in\relint Q.
\]
Suppose \(y\in F_i\).  By
\eqref{eq:low-dimensional-facet-image}, \(y\notin E\) forces
\(\dim F_i=j-1\).  Thus \(F_i\) is a \((j-1)\)-simplex and
\[
 \partial F_i\subseteq f((\Delta^j)^{(j-2)})=E.
\]
Hence \(y\in\relint F_i\).
Lemma~\ref{lem:positive-cone-section}, applied with
\(I=\{0,\ldots,j\}\setminus\{i\}\), gives
\(z=\kappa_\ell(y)\in\relint\kappa_\ell(F_i)\).
Since \(\mathcal K\) subdivides
\(\kappa_\ell(F_i)\), Lemma~\ref{lem:refinement-membership}(i) gives
\[
 Q\subseteq\kappa_\ell(F_i)
 \quad\Longleftrightarrow\quad
 y\in F_i.
\]
The reverse implication in this display is the lemma; the forward
implication is immediate from \(y=\kappa_\ell^{-1}(z)\).

It follows from Lemma~\ref{lem:generic-degenerate-fiber} that
\[
 \mu_{\sum_i\widehat P_i}(Q)
 =
 \#\{i:y\in F_i,\ \dim F_i=j-1\}=0
 \quad\text{in }\F.
\]
The coefficient criterion gives
\[
 \sum_iP_i=\varpi_{j-1}\!\left(\sum_i\widehat P_i\right)=0,
\]
which is the desired cancellation.
\end{proof}

\end{document}
